%%%%%
%
%% Dans ce fichier, les �ventuels accents sont cod�s
%% en ISO-8859-15 (aka latin9)
%% La ligne suivante donne trois exemples 
%% e accent aigu : �, e accent grave : �, e-dans-l'o : �
\section{Estimation des hauteurs et bases de nuage � partir d'un profil lidar}

\itemstyle{Op\'erateur~:} \operfmt{Nuage\_lidar}
\index{Nuage\_lidar}

%%%%%OBJET%%%%
\itemstyle{Objet~:}
Cet op�rateur d�termine la hauteur des bases et sommets des nuages par une m�thode mise au point par A. Mortier 2013 \cite{THESE-MORTIER-2013} et inspir�e de Pal et al. (1992). 

%%%%%ARGUMENTS%%%%
\bigskip
%% Dans ce fichier, les �ventuels accents sont cod�s
%% en ISO-8859-15 (aka latin9)
%% La ligne suivante donne trois exemples 
%% e accent aigu : �, e accent grave : �, e-dans-l'o : �
%%%%%ENTREES%%%%
\begin{tabarg}
\multicolumn{3}{l}{\titretab {Entr\'ees}}\\
\hline
$X$ & Vecteur, \texttt{float[$N_{level}$]} & Signal lidar \\
$S_1$ & Coeff., \texttt{int} & Seuil de d�tection \\
$S_2$ & Coeff., \texttt{int} & Seuil de d�tection \\
$Z$ & Vecteur, \texttt{float[$N_{level}$]} & Altitudes r�elles des portes du lidar (m)\\
$Z_{min}$ & Coeff., \texttt{int} & Altitude en dessous de laquelle on ne d�tecte pas de nuage (m)\\
\end{tabarg}

\medskip
%%%%%SORTIES%%%%
\begin{tabarg}
\multicolumn{3}{l}{\titretab {Sorties}}\\
\hline
$Base$ & Vecteur, \texttt{float[6]} & Altitude des bases des nuages (m)\\
$Max$ & Vecteur, \texttt{float[6]} & Altitude des maximums de r�tro-diffusion (m)\\
$Sommet$ & Vecteur, \texttt{float[6]} & Altitude des sommets des nuages (m)\\
$Flag_{nuage}$ & Vecteur, \texttt{float} & Nombre de couches d�tect�es (m)\\
\end{tabarg}



%%%%%FORMULE%%%%
\itemstyle{Formule}

%% Formule de l'operateur : Nuage_lidar
%% Dans ce fichier, les �ventuels accents sont cod�s
%% en ISO-8859-15 (aka latin9)
%% La ligne suivante donne trois exemples 
%% e accent aigu : �, e accent grave : �, e-dans-l'o : �


L'algorithme est d�crit dans la th�se cit�e en r�f�rence dans la partie 2.3.4.2. Il est bas� sur la variabilit� verticale de chaque profil. Les minimuns et maximums du profil sont dans un premier temps d�termin�s � partir de la d�riv�e du profil. Puis ces valeurs sont s�lectionn�es en fonction des 2 seuils $S_1$ et $S_2$.\\
Pour qu'un nuage soit identifi� il faut que : \\
\begin{itemize}
\item la diff�rence du signal lidar entre le maximum et la base d�tect�s soit sup�rieur � $S_1$ fois le bruit du signal 
\item le rapport de signal lidar entre le maximum et le sommet du nuage d�tect� soit sup�rieur � $S_2$.
\end{itemize}
Les valeurs recommand�es pour les seuils de d�tection sont $S_1= 10$ et $S_2=4$.

 

%%%%%AUTEUR%%%%
\itemstyle{Auteur~:}
CNRM/GMEI/TRAMM

%%%%%REFERENCE%%%%
\itemstyle{R\'ef\'erences~:}
Utilis� dans la fili�re BASIC du LOA, A. Mortier 2013 \cite{THESE-MORTIER-2013}.
